% ============================================================
% f2014_sichuan_arts_q06_flowchart.tex
% 2014 四川卷（文科）第 6 题：程序框图
% 来源 PDF：普通高考/2014/2014四川文.pdf
% 原页：第 1 页（A3 横排）中栏上方，印面 921，第 6 题题干右侧
% 结构：开始 → 输入 x,y → 判断「x≥0，y≥0，x+y≤1？」；
%       是：S=2x+y；否：S=1；两路汇合 → 输出 S → 结束。
% ============================================================
\begin{tikzpicture}[>=Stealth, line width=0.7pt, font=\small,
  x=1cm, y=1cm, line cap=round, line join=round]

  % 开始
  \draw[rounded corners=6pt] (-0.95,0.45) rectangle (0.95,-0.45);
  \node at (0,0) {开始};
  \draw[->] (0,-0.45) -- (0,-0.90);

  % 输入
  \draw (-1.20,-0.90) -- (1.70,-0.90) -- (1.20,-1.60) -- (-1.70,-1.60) -- cycle;
  \node at (0,-1.25) {输入 $x$，$y$};
  \draw[->] (0,-1.60) -- (0,-2.00);

  % 判断（六边形，与原页一致）
  \draw (-2.90,-2.90) -- (-2.30,-2.40) -- (2.30,-2.40) -- (2.90,-2.90)
     -- (2.30,-3.40) -- (-2.30,-3.40) -- cycle;
  \node at (0,-2.78) {$x\geqslant 0$，$y\geqslant 0$，};
  \node at (0,-3.12) {$x+y\leqslant 1$？};

  % 是：左支
  \draw[->] (-2.90,-2.90) -- (-3.90,-2.90) -- (-3.90,-3.70);
  \node at (-3.15,-2.62) {是};
  \draw (-5.20,-3.70) rectangle (-2.60,-4.50);
  \node at (-3.90,-4.10) {$S=2x+y$};

  % 否：右支
  \draw[->] (2.90,-2.90) -- (3.90,-2.90) -- (3.90,-3.70);
  \node at (3.15,-2.62) {否};
  \draw (2.60,-3.70) rectangle (5.20,-4.50);
  \node at (3.90,-4.10) {$S=1$};

  % 两路汇合
  \draw (-3.90,-4.50) -- (-3.90,-4.90) -- (3.90,-4.90) -- (3.90,-4.50);
  \draw[->] (0,-4.90) -- (0,-5.25);

  % 输出
  \draw (-1.10,-5.25) -- (1.80,-5.25) -- (1.30,-5.95) -- (-1.60,-5.95) -- cycle;
  \node at (0.10,-5.60) {输出 $S$};
  \draw[->] (0.10,-5.95) -- (0.10,-6.40);

  % 结束
  \draw[rounded corners=6pt] (-0.85,-6.40) rectangle (1.05,-7.30);
  \node at (0.10,-6.85) {结束};
\end{tikzpicture}
